O algoritmo Selection Sort apresenta comportamento semelhante nos
diferentes tipos de entrada, pois a quantidade de comparações realizadas
não depende da organização inicial dos elementos do vetor. Sua complexidade
de tempo é analisada considerando o melhor caso, o caso médio e o pior caso.

\subsubsection{Melhor caso -- \texorpdfstring{$\Theta(n^2)$}{Theta(n²)}}

O melhor caso ocorre quando o vetor já está ordenado de forma crescente.
Mesmo nessa situação, o algoritmo precisa percorrer a parte ainda não
ordenada do vetor para verificar qual é o menor elemento.

Considerando um vetor com $n$ elementos, na primeira iteração são
realizadas $n-1$ comparações. Na segunda iteração, são realizadas $n-2$
comparações, e assim sucessivamente até a última iteração, que realiza
apenas uma comparação.

A quantidade total de comparações pode ser representada pela soma:

\begin{equation}
T(n) = (n-1) + (n-2) + \cdots + 2 + 1
\end{equation}

Essa soma corresponde à soma dos primeiros $n-1$ números naturais:

\begin{equation}
T(n) = \frac{n(n-1)}{2}
\end{equation}

Desenvolvendo a expressão:

\begin{equation}
T(n) = \frac{n^2-n}{2}
\end{equation}

O termo de maior crescimento é $n^2$. Portanto:

\begin{equation}
T(n) = \Theta(n^2)
\end{equation}

Assim, o melhor caso do Selection Sort apresenta complexidade
quadrática.

\subsubsection{Pior caso -- $\Theta(n^2)$}

O pior caso ocorre quando o vetor está organizado de forma que seja
necessário realizar várias trocas para colocá-lo em ordem crescente.
Entretanto, mesmo nesse cenário, a quantidade de comparações realizadas
pelo algoritmo permanece a mesma.

Na primeira iteração, são realizadas $n-1$ comparações; na segunda,
$n-2$ comparações, e assim sucessivamente. Dessa forma:

\begin{equation}
T(n) = (n-1) + (n-2) + \cdots + 2 + 1
\end{equation}

Aplicando a fórmula da soma:

\begin{equation}
T(n) = \frac{n(n-1)}{2}
\end{equation}

\begin{equation}
T(n) = \frac{n^2-n}{2}
\end{equation}

Como o termo dominante é $n^2$, temos:

\begin{equation}
T(n) = \Theta(n^2)
\end{equation}

Portanto, o pior caso do Selection Sort apresenta complexidade
quadrática.

\subsubsection{Caso médio -- $\Theta(n^2)$}

O caso médio ocorre quando os elementos estão distribuídos de maneira
aleatória. Nessa situação, assim como nos demais casos, o algoritmo
precisa percorrer toda a parte ainda não ordenada do vetor para encontrar
o menor elemento.

A quantidade de comparações continua sendo representada por:

\begin{equation}
T(n) = (n-1) + (n-2) + \cdots + 2 + 1
\end{equation}

Aplicando a fórmula da soma:

\begin{equation}
T(n) = \frac{n(n-1)}{2}
\end{equation}

\begin{equation}
T(n) = \frac{n^2-n}{2}
\end{equation}

Como o termo dominante é $n^2$, a complexidade assintótica é:

\begin{equation}
T(n) = \Theta(n^2)
\end{equation}

Portanto, no caso médio, o Selection Sort também apresenta complexidade
quadrática.

\subsubsection{Resumo das complexidades}

\begin{table}[H]
    \centering
    \caption{Complexidades do algoritmo Selection Sort}
    \label{tab:complexidade-bubble}
    \begin{tabular}{lc}
        \hline
        \textbf{Caso} & \textbf{Complexidade de tempo} \\
        \hline
        Melhor caso & $\Theta(n^2)$ \\
        Caso médio & $\Theta(n^2)$ \\
        Pior caso & $\Theta(n^2)$ \\
        \hline
    \end{tabular}
\end{table}

Conclui-se que o Selection Sort apresenta complexidade quadrática nos
três casos analisados. Diferentemente do Insertion Sort, a organização
inicial do vetor não reduz significativamente a quantidade de comparações,
pois o algoritmo precisa procurar o menor elemento em toda a parte ainda
não ordenada a cada iteração.
